% Hand-written sentences that depend on results; every number is a macro from v19_numbers.tex (Research/code/v19_tables.py).
\newcommand{\nPointerSentence}{by its certificates, fixed forms lose $4.00$ and $8.00$, the maximum that Theorem~\ref{thm:universal} allows, two rounds lose $1.80$ and between $3.32$ and $3.33$, against the upper bounds $2$ and $4$ of part (a), and $b-1$ rounds lose $1.80$ at both caps (Appendix~\ref{app:pointerexact}).}
\newcommand{\nOrderSentence}{an illustrative flip of the last answer with probability 10\% biases the scores by \nPosBranch--\nPosFit{} standard errors at $n=1000$ (Appendix~\ref{app:measurement}).}
\newcommand{\nEmAppendixText}{At $n=1000$ the working model on the fitted-form records has the smallest RMSE of the four estimators in most settings, \nEmBrLo--\nEmBrHi{} times that of branching: the model shares strength across paths that the design-based scores treat separately. Its bias is between \nEmBiasLo{} and \nEmBiasHi{} branching standard deviations and its coverage \nEmCovLo--\nEmCovHi{} (below 0.925 in \nEmBelow{} of \nEmSettings{} settings). The bias is a property of the model, not of the sample size: the accuracy of a depth-$b$ tree depends on interactions of up to $b$ items and the label, which a two-way model does not represent. At $n=\nEmBigN$ the same bias is up to \nEmBigBiasHi{} of the estimator\textquotesingle s own standard deviations, coverage falls to \nEmBigCovLo--\nEmBigCovHi{} (below 0.90 in \nEmBigBelow{} of \nEmBigSettings{} settings), and the RMSE is \nEmBigRmseLo--\nEmBigRmseHi{} times that of branching and above that of the design-based form score in \nEmBigWorse{} settings; the design-based scores keep coverage \nEmBigFixCovLo--\nEmBigFixCovHi. A richer model would reduce the bias and raise the variance; with all interactions it is the design-based estimator. On balanced random forms the model-based estimate is unbiased to a similar degree but \nEmUnifLo--\nEmUnifHi{} times less accurate at $n=1000$ than the design-based score on optimized forms, so the choice of forms matters for model-based analysis as well.}
\newcommand{\nDiscussionSentence}{On the trained rules we examined, with the cap equal to the rule's depth, optimal forms needed between $1.0$ and \nMaxTrees{} times the records of branching where the optimum was computed; one extra question lowered the median from \nAugZeroMed{} to \nAugOneMed, and one adaptive round came within \nTwoBoundMax{} of branching in every screening run; per question, the best of the three pre-committed designs we priced cost \nQBest{} times branching at the median.}
\newcommand{\nAbstractData}{On screening trees for several health outcomes, with the cap equal to the depth of the rule, optimal forms need $1.0$--\nMaxTrees{} times the records of the adaptive optimum where the optimum was computed (at most \nPsBigHi{} for the widest rules), and \nSmoothMedLo--\nSmoothMedHi{} at the median when the design is learned from a pilot charged to the same audit. Trained rules depend on few items beyond their depth, so one extra question per record lowers the median price from \nAugZeroMed{} to \nAugOneMed{} and two to \nAugTwoMed; counted in questions instead of records, however, branching keeps an advantage of about \nQBest{} over the best of the pre-committed designs we priced (forms of size $b$ and $b+1$, and complete records of the rule's items).}
\newcommand{\nContribData}{On screening trees for several health outcomes, with the cap equal to the depth of the rule, optimal forms need $1.0$--\nMaxTrees{} times the records of branching, on the rules whose optimum was computed. This is far from the worst case, and in records it is easy to remove: the rules depend on a median of \nExcessMed{} items beyond their depth, and one extra pre-committed question lowers the median price from \nAugZeroMed{} to \nAugOneMed, two to \nAugTwoMed, as does one follow-up round. In questions it is not removed: forms of size $b$, forms one question longer and complete records of the rule's items cost \nQForms, \nQFormsOne{} and \nQAll{} times the questions of branching at the median. Uniformly drawn forms cost \nUitMedLo--\nUitMedHi{} times the records of branching, so what matters first is optimizing the forms, which the solver does from a reference sample or a charged pilot.}
\newcommand{\nAuditParagraph}{In simulated audits of 500 to 32,000 respondents on \nCurveTrees{} trees (Figure~\ref{fig:rounds}c, Table~\ref{tab:curve}), the mean squared error of forms relative to branching is \nAudDesignLo--\nAudDesignHi{} for the design-based scores and \nAudTwoFormsSurveyLo--\nAudTwoFormsSurveyHi{} for the pairwise EM estimate up to 8,000 respondents; at 32,000 the bias shows, the pairwise estimate is less accurate than the design-based form score for \nCurveTwoWorsed{} of the \nCurveTrees{} trees, and its intervals cover \nEmBigCovLo--\nEmBigCovHi{} (Table~\ref{tab:em}).}
\newcommand{\nAuditAppendixText}{For the design-based scores the ratio of mean squared errors stays near the population ratio at every size (medians \nAudDesignLo--\nAudDesignHi). For the three-way EM estimate the ratio between forms and branching records is \nAudThreeFormsLo--\nAudThreeFormsHi{} on the fitted forms and \nAudThreeMformsLo--\nAudThreeMformsHi{} on the forms optimized for the pairwise model; for the pairwise model it is \nAudTwoFormsSurveyLo--\nAudTwoFormsSurveyHi{} and \nAudTwoMformsSurveyLo--\nAudTwoMformsSurveyHi{} up to 8,000 respondents (at most \nAudTwoFormsSurveyMax{} and \nAudTwoMformsSurveyMax{} for a single tree). These orders of magnitude agree with the population constants of Table~\ref{tab:model}, although the target laws do not satisfy the models. Branching here is the rule's own tree, not the branching design that is optimal within the model, which Table~\ref{tab:model} puts a few percent below the optimal forms. At large sizes the comparison is governed by the bias of the model, on either kind of record: at 32,000 respondents the pairwise fit has a median absolute bias of \nAudBiasFormsBig{} of its standard deviation on forms (\nAudBiasFormsSmall{} at 1,000) and \nAudBiasBranchBig{} on branching records (\nAudBiasBranchSmall{} at 1,000), the pairwise fit to branching records is less accurate than the plain mean of $t$ for \nAudBranchTwoWorsec{} and \nAudBranchTwoWorsed{} of the \nCurveTrees{} trees at 16,000 and 32,000, and the ratio between the two designs ranges from \nAudTwoFormsMinAll{} to \nAudTwoFormsMaxAll{} over trees. Against the design-based form score on the same fitted-form records, the pairwise EM estimate has the smaller mean squared error for every tree up to $n=\nCurveLastAllBetter$ and the larger for \nCurveTwoWorsec{} and \nCurveTwoWorsed{} of the \nCurveTrees{} trees at 16,000 and 32,000 (RMSE up to \nCurveTwoMaxc{} and \nCurveTwoMaxd{} times larger); which trees these are cannot be read from the capped records. A model-assisted score, which uses a working model only to choose the tables of a pointwise-unbiased score, as in the smoothed pilot of Appendix~\ref{app:pilot}, cannot have a smaller variance than the oracle tables for the same form distribution; the advantage of the model-based estimate at small $n$ is bought with the model's bias.}
\newcommand{\nForestScaleSentence}{On twenty Spambase features the picture is the same at $b=3$: forms need \nFScaleThreeLo--\nFScaleThreeHi{} times the records of branching and two rounds at most \nFScaleThreeTwoHi{} (Table~\ref{tab:forestscale}).}
\newcommand{\nForestScaleAppendixText}{At $b=3$ column generation over branching interviews closed within its budget on all three splits, so the ratios \nFScaleThreeLo--\nFScaleThreeHi{} are exact to the certified gap; they lie in the range of the twelve-feature tasks. The two-round programs did not close in 15 minutes, and their certified intervals put two rounds between \nFScaleThreeTwoLo{} and \nFScaleThreeTwoHi{} times the records of branching. At $b=4$ (4,845 forms) the fixed optimum was certified in seven minutes, but 25 minutes of column generation over branching interviews left only the trivial lower bound for the one split that was run, so its interval is uninformative; the other two splits were not run. Exact adaptive designs at $m=20$ and $b\ge4$ are beyond this implementation.}
\newcommand{\nDependenceAppendixText}{Along the strongly coupled path the pairwise constants of both classes rise from $1/3$ towards their values without a model, and their ratio rises from $1.00$ at $\beta=0$ to \nDepRatioLast{} at $\beta=\nDepBetaMax$, above the limit $2^{b-1}/b=1.33$ guaranteed by the theorem and below the factor 4 of part (a). At moderate coupling the gap is still small (\nDepRatioOne{} at $\beta=1$ and \nDepRatioTwo{} at $\beta=2$); it opens where the auxiliary features become nearly deterministic functions of the address and the data. At the pairwise fits of the noisy address populations the ratio stays at \nDepFitTwoMax{} at every accuracy, although the ratio without a model is $2^{b-1}$ at each of these populations: the association of the label with the data bits, which is what accuracy adds, does not bring the address structure into the pairwise tangent space. At the three-way fits the ratio is \nDepFitThreeLo--\nDepFitThreeHi.}
\newcommand{\nAipwAppendixText}{The cross-fitted estimator has the mean squared error of the reference-table score to within a few percent (median ratio \nAipwMseLo--\nAipwMseHi{} over the three sizes, at most \nAipwMseMax{} for a single tree), a bias within Monte Carlo error ($|$bias$|$/SD at most \nAipwBiasMax, against \nAipwBiasRefMax{} for the reference-table score on the same records), and coverage between \nAipwCovLo{} and \nAipwCovHi{} over trees and sizes (\nAipwCovRefLo--\nAipwCovRefHi{} for the reference-table score; the low values are for the same tree and size). A design-unbiased analysis therefore needs no external reference. It does not capture the small-sample advantage of the pairwise EM estimate, as expected: for given form probabilities its variance cannot fall below that of the oracle tables, and the advantage of the model-based estimate is bought with the model's bias.}
\newcommand{\nPracticeSentence}{If the rule uses at most $b$ features, ask them all; if the interview can branch, administer the rule and take a plain mean. If questions must be pre-committed, do not randomize uniformly: solve the path-set program on a reference sample or on a pilot that is charged to the audit, draw one form per record from the resulting distribution, and average the scores, which are unbiased whatever the reference was. If records are scarce and the burden allows one or two questions more than the rule asks, or one batch of follow-ups, use them; if questions are scarce, no pre-committed design we priced helps, and only adaptivity, full branching or one follow-up round, saves them. A pairwise working model on pair-covering forms is nearly as efficient as branching for audits of a few thousand records, but its bias does not shrink, so claims that must not rest on a model should use the design-based scores.}
\newcommand{\nRulesSentence}{Trees for four other BRFSS outcomes (21 variables) and for Spambase (57 features), with forms on the rule's items, give \nRulesLo--\nRulesHi{} in the \nRulesOk{} runs with a gain that were computed (median \nRulesMed; Table~\ref{tab:rules}). For the \nRulesSkipped{} rules with more than 13 items the optimum over all forms was not computed; the optimized design on the tree's own path sets, a convex program whose size does not grow with the number of items (Proposition~\ref{prop:upper}), certifies ratios of at most \nPsBigLo--\nPsBigHi{} for them, and it is within \nPsAllMax\% of the optimum on the \nPsAllN{} rules where both are known (median \nPsAllMed\%). Bootstrap resamples of the target records move a ratio by at most $\nRulesBootBrfss\%$ on BRFSS and $\nRulesBootSpam\%$ on the \nRulesSpamN{} Spambase target records.}
\newcommand{\nMisspecSentence}{The accuracy of a depth-$b$ tree involves interactions of up to $b$ items and the label, so the target laws do not satisfy the pairwise model, and the comparison is one of efficiency against robustness. At the target laws themselves we compute the limit of the pairwise maximum-likelihood estimate under each design and its sandwich variance \citep{white1982,golden2019} (\nMisRuns{} runs; Table~\ref{tab:misspec}). The estimate is well defined essentially when the design observes every pair of items together, the classical condition for split questionnaires \citep{raghunathan1995,rhemtulla2012}; the rule's own tree and the optimal design-based forms mostly do not. On pair-covering forms optimized for the model, the limit misses the true accuracy by at most \nMisBiasPtsPc{} points, the squared bias equals the variance at a median of \nMisCrossPc{} records, and the sandwich variance is \nMisVarPmLo--\nMisVarPmHi{} times that of the branching design that is optimal for the model (median \nMisVarPmMed).}
\newcommand{\nRulesAppendixText}{Of the \nRulesAll{} runs, \nRulesOk{} have a gain and were computed, \nRulesTrivial{} are trivial, and \nRulesSkipped{} were not computed (\nRulesTooMany{} rules depend on more than 13 items; \nRulesUnfinished{} did not finish). For \nRulesFullN{} of the computed runs the bootstrap, the uniform design and the pairwise fit were computed as well; for the other \nRulesFilled{} only the optimal forms. The ratios lie in the range of the diabetes trees of Table~\ref{tab:screening}, with a few larger values (up to \nRulesHi). For the runs that were not computed, the lower end of the bracket is \nRulesBrkLo--\nRulesBrkHi, within the range of the computed ratios, and the upper end \nRulesBrkUpLo--\nRulesBrkUpHi{} is not informative; the largest rules are thus not known to behave differently, but their ratios are not known either. Resampling the target records with the tree fixed moves a ratio by at most \nRulesBootBrfss\% of its value on BRFSS, where a target half has more than $10^5$ records, and by at most \nRulesBootSpam\% on Spambase (\nRulesSpamN{} target records; 5th to 95th percentile of \nRulesBootReps{} resamples, \nRulesBootN{} runs in all). The variation in the tables is therefore between rules and training samples, not sampling error of the target law; the bootstrap does not cover the uncertainty of training the rule. Uniform randomization over the forms of the rule\textquotesingle s items needs \nRulesUnifLo--\nRulesUnifHi{} times the records of branching (median \nRulesUnifMed), and the ratio within the pairwise model is \nRulesModelLo--\nRulesModelHi{} (median \nRulesModelMed, \nRulesModelN{} runs).}
\newcommand{\nMisspecAppendixText}{Four findings are stable over the \nMisRuns{} runs. First, identification follows pair coverage. MF, PC and MB observe every pair of the rule\textquotesingle s items together in every run; F does in \nMisCovF{} runs and T in \nMisCovT. Over the designs and runs that cover all pairs, all \nMisStarts{} starts reach one interior maximizer in \nMisCovWell{} of \nMisCovCells{} cases and none is found in \nMisCovNone; over those that leave a pair uncovered, the counts are \nMisUncWell{} and \nMisUncNone{} of \nMisUncCells. Coverage with a small probability is not always enough in practice: MF, whose smallest pair probability can be below 1\%, has one maximizer from all starts in \nMisWellMF{} runs and MB in \nMisWellMB, whereas PC, with its 10\% floor, has one in all \nMisWellPC. Second, on PC the bias is at most \nMisBiasSdPc\% of the standard deviation of $t$, or \nMisBiasPtsPc{} accuracy points, the sandwich variance is between \nMisSandPcLo{} and \nMisSandPcHi{} of $\Var_P(t)$, below the variance of the design-based mean on the tree, and the floor costs little: the variance of PC is \nMisVarPfLo--\nMisVarPfHi{} times that of MF where both are defined. Third, where PC and MB both have an interior maximizer (\nMisBoth{} runs), the ratio of their sandwich variances is \nMisVarPmLo--\nMisVarPmHi{} (median \nMisVarPmMed), MB, optimal for the fitted law, has the larger variance at the true law in \nMisMbWorse{} runs, and with the bias included the mean squared error of PC relative to MB is a median \nMisMsePma, \nMisMsePmb{} and \nMisMsePmc{} at 1,000, 4,000 and 16,000 records (ranges \nMisMsePmLoa--\nMisMsePmHia{} and \nMisMsePmLoc--\nMisMsePmHic{} at the two ends). Fourth, the variance computed from the model\textquotesingle s information matrix is within a factor \nMisSandModelPc{} of the sandwich variance on PC, but up to \nMisSandModelHi{} times too small on other designs, which contributes to the under-coverage of model-based intervals in Table~\ref{tab:em} beyond the effect of bias. The design-based ratio at the same laws is \nMisDesignLo--\nMisDesignHi{} (median \nMisDesignMed). The asymptotic bias agrees with the simulated audits: for the trees of caps four and five the calculation puts the bias of PC at a median \nMisBiasBig{} standard deviations at 32,000 records, and the EM estimate on the model forms of the audits, which have a 5\% floor, has \nAudBiasMformsBig.}
\newcommand{\nCiSentence}{In the model of Theorem~\ref{thm:structured}(b), with features independent given the label, the constants are in closed form for every rule, including those whose design-based optimum was not computed: at the fit, forms drawn without the label need at most \nCiBrfssPct\% more records than the branching optimum on every BRFSS rule and at most \nCiSpamPct\% on Spambase (\nCiRuns{} rules; Table~\ref{tab:ci}).}
\newcommand{\nCiAppendixText}{The closed forms of Theorem~\ref{thm:structuredfull}(b) need only the class-conditional degree-one coefficients of the rule, so they are available for every rule: the \nCiRuns{} rules of Tables~\ref{tab:screening} and~\ref{tab:rules} that depend on more than $b$ items. At the conditional-independence fit of the target law, forms drawn without the label need a median \nCiMedBrfss\% and at most \nCiBrfssPct\% more records than the optimum with label-dependent forms on the BRFSS rules, and at most \nCiSpamPct\% on Spambase, where the classes differ more; the branching optimum without the label lies between the two. Under this working model, then, pre-commitment costs nothing that matters on any rule we trained, at caps four to six and for rules with up to 16 items. The model is cruder than the pairwise model: it ignores the associations among the items, which the target laws have. Its constants for forms are \nCiDiabLo--\nCiDiabHi{} of the variance of $t$ on the diabetes trees (\nCiBlindLo--\nCiBlindHi{} over all rules), and its bias for the true accuracy was not computed; we report it as the case in which the value of adaptivity is known in closed form, not as a recommended analysis.}
\newcommand{\nAugSentence}{With forms longer than the rule is deep, the price in records falls much faster than in the worst case of Theorem~\ref{thm:rounds}(c). Over the \nAugOneN{} rules with a gain for which both were computed, optimal forms need a median \nAugZeroMed{} times the records of branching at size $b$, \nAugOneMed{} at size $b+1$ (at most \nAugOneMax) and \nAugTwoMed{} at size $b+2$ (at most \nAugTwoMax; Figure~\ref{fig:rounds}b, Table~\ref{tab:capcurve}). The reason is the width of trained rules, not a failure of the theorem: these trees depend on a median of \nExcessMed{} and at most \nExcessMax{} items beyond their depth, whereas the address contrast depends on $2^{b-1}-1$. \emph{Counting questions.} Extra questions save records and hardly any questions. Per question asked, and relative to a branching interview charged $b$ questions per record, forms of size $b$ cost a median \nQForms, forms of size $b+1$ cost \nQFormsOne, and asking all $K$ items of the rule in a fraction $b/K$ of the records, a mean-budget design, costs \nQAll{} (\nQAllLo--\nQAllHi). Against each of these three pre-committed designs branching saves between a quarter and a third of the questions at the median (the best of the three, rule by rule, costs \nQBest); a two-round interview with $b$ questions saves them as well. Which design is best depends on whether records or questions are the scarce resource.}
\newcommand{\nCapcurveAppendixText}{Of the \nAugRules{} rules that depend on more than $b$ items and were run (\nAugMissing{} further rule, with eleven items at cap six, was stopped unfinished), \nAugOneN{} have the optima at sizes $b$ and $b+1$ and \nAugTwoN{} also at $b+2$; the others have more items than the solver handles at those sizes. At size $b+1$ the forms need a median \nAugOneMed{} times the records of branching (at most \nAugOneMax), against \nAugZeroMed{} at size $b$: a median \nAugOneRec\% and at least \nAugOneRecMin\% of the log-gap is removed. At size $b+2$ the median is \nAugTwoMed{} (at most \nAugTwoMax; \nAugTwoRec\% removed in the median). Theorem~\ref{thm:augment} allows much less in the worst case, a factor two for one extra question at any cap. The difference is the number of items beyond the depth, $K-b$, which is \nExcessMed{} in the median and at most \nExcessMax{} for the rules computed at size $b+1$ (at most \nExcessAllMax{} over all rules) and $2^{b-1}-1$ for the address contrast. The comparison is with the branching interview of depth $b$, which already attains the full-information variance for a single tree, so extra questions cannot help branching. The optimized path-set design of size $b$ needs a median \nPsAllMed\% and at most \nPsAllMax\% more records than the optimal forms on the \nPsAllN{} rules where both were computed; for the \nPsBigN{} rules whose optimum at size $b$ was not computed it certifies that forms need at most \nPsBigLo--\nPsBigHi{} times the records of branching, against the lower ends \nRulesBrkLo--\nRulesBrkHi{} of the bracket. The values in this table are primal values of the program over all forms of the given size, accurate to the tolerance of the cone solver (about $10^{-3}$ in relative terms).}
